\begin{afterword}
\summaryhead

A short story, the last in the book set among the beisutsukai. At dawn on Mount Mirror, Jeffreyssai tells his five students that their training is over. He cannot make masters. What he taught is incomplete and will hurt them, and mastery comes from using it until it fails. Asked why he taught them at all, he answers that a taught student who fails may recover and remake the Art, where an untaught one might think Reason itself had failed. He does not know whether this is the only road. He warns of three flaws: looking harder for flaws in arguments one would rather not accept, cleverness, and underconfidence that feels like humility. Then he leaves, saying that one of them seems likely to come back.

Alone, Brennan finds that he does not know what he wants, and that his pursuit of power may have been a role. Asking why he might not know a strong desire, he sees that he has hidden one from himself because it seems impossible. He sets out to pursue it, realizes that Jeffreyssai foresaw this, and rejects the proverb that it is futile to ask how a master knows things.

\responsehead

The story is fiction and is judged for what it shows; the teacher's points are judged as arguments.

The story is well built, because its second half acts out its first. Brennan falls into the third flaw on the same day, when he applies the doubt he was taught to his picture of himself. He gets out by a question his training gave him, as Jeffreyssai told Yin a taught student might. The story itself suggests an answer to the proverb Brennan rejects at the end. Jeffreyssai read a flicker of Brennan's eyelids, and Taji's stated wish was weak evidence of what Taji wants, so a master's knowing looks like inference from small signs. And Brennan's refusal to let the proverb stop his curiosity is the point of ``Science as Curiosity-Stopper,'' early in the book.

The teacher's claims come as a teacher's experience, and the story keeps his hedges. He does not know whether failure is the only road, and he says an untaught student's chances would be ``less,'' not nil. One claim is put more firmly: when the Art seems to fail, ``the reality will be that it was you who failed your Art.'' Stated as fact, this places every failure on the student and none on the Art, and the story gives no way to tell the two apart in a given case. As advice on how to take a failure, it is the attitude that ``Something to Protect'' recommends.

The three flaws gather the book's earlier warnings. The first restates ``Knowing About Biases Can Hurt People.'' The third condenses ``The Sin of Underconfidence,'' posted a week before this story and placed later in the book, which named underconfidence as the third of three sins without naming the other two. The second comes with its own account of why a known weakness persists: the clever value cleverness for itself.

What the story leaves out is the part that would test its central claim. The failure through which mastery is said to come is foretold, not shown; what the reader sees is the teaching helping in a small way on the first day. Brennan's goal is not named. Like ``The Ritual'' and ``Class Project,'' the story stops before any result.

In short: a well-made farewell whose second half acts out its first and suggests an answer to its own closing question; like ``The Ritual'' and ``Class Project,'' it stops before the failure through which, it says, mastery comes.
\end{afterword}
